% =============================================================================
% CHƯƠNG 1 — GIỚI THIỆU ĐỀ TÀI, CÁC BÊN LIÊN QUAN VÀ YÊU CẦU NGHIỆP VỤ
% Nội dung đối chiếu với repo marketpulse-360 @ dbe614b.
% =============================================================================
\setcounter{page}{1}

\section{Bối cảnh và bài toán}
\label{sec:c1-boi-canh}

\subsection{Bài toán marketing đa kênh trong bán lẻ}
\label{sec:c1-bai-toan}

Các doanh nghiệp bán lẻ tại Việt Nam hiện nay hiếm khi chỉ bán hàng qua một kênh. Cùng một lúc, doanh nghiệp chạy nhiều chiến dịch quảng cáo trên nhiều nền tảng khác nhau (Google Ads, Meta Ads, TikTok Ads và Shopee Ads), đưa người dùng về website và ứng dụng di động, ghi nhận đơn hàng trên hệ thống thương mại điện tử, quản lý hồ sơ khách hàng trong CRM và nuôi dưỡng khách hàng bằng các chiến dịch email. Mỗi hệ thống trên tự sinh ra báo cáo riêng với định dạng, độ trễ, đơn vị tiền tệ và định danh khác nhau: báo cáo quảng cáo của Google tính chi phí bằng \texttt{cost\_micros}, của Meta tính bằng chuỗi thập phân, của Shopee ghi ngày theo dạng \texttt{DD/MM/YYYY}; sự kiện web đến theo thời gian thực trong khi báo cáo quảng cáo và ảnh chụp CRM đến theo ngày.

Hệ quả là những câu hỏi kinh doanh cơ bản nhất, như \emph{"chiến dịch nào thực sự mang lại doanh thu?"}, \emph{"ngân sách hôm nay có đang bị đốt vào kênh kém hiệu quả không?"} hay \emph{"khách hàng rời bỏ ở bước nào của phễu?"}, lại không thể trả lời bằng bất kỳ báo cáo đơn lẻ nào. Các nền tảng quảng cáo cùng ``nhận công'' cho một đơn hàng, số liệu đơn hàng không khớp với số liệu chuyển đổi do nền tảng báo cáo, và mỗi phòng ban làm việc với một con số khác nhau.

Nhằm giải quyết vấn đề này, nhóm xây dựng đề tài \textbf{MarketPulse 360}, một sản phẩm dữ liệu marketing đầu-cuối nhằm chứng minh rằng giải pháp Kỹ thuật Dữ liệu có thể vừa \emph{đúng}, vừa \emph{nhanh}, vừa \emph{hữu ích} cho người dùng nghiệp vụ. Hệ thống tiếp nhận năm nhóm nguồn dữ liệu, xử lý song song theo hai nhánh luồng và lô, kiểm chứng tự động trước khi công bố, rồi phục vụ kết quả cho ứng dụng \emph{Marketing Command Center} (Hình~\ref{fig:system-context}). Toàn bộ mã nguồn, dữ liệu mô phỏng và bằng chứng kiểm thử được lưu trong kho mã của nhóm~\cite{mp360repo}.

\begin{figure}[H]
    \centering
    \resizebox{\textwidth}{!}{% Hình 1.1 — Bối cảnh hệ thống và luồng giá trị của MarketPulse 360
\begin{tikzpicture}[node distance=2.5mm]
  % Cột 1: nguồn dữ liệu
  \node[mpbox=40mm] (ads) {\textbf{Báo cáo quảng cáo}\\ Google, Meta,\\ TikTok, Shopee};
  \node[mpbox=40mm, below=of ads] (web) {\textbf{Sự kiện web/app}\\ hành vi, phiên, click};
  \node[mpbox=40mm, below=of web] (orders) {\textbf{Đơn hàng}\\ tạo đơn, hoàn tiền};
  \node[mpbox=40mm, below=of orders] (crm) {\textbf{CRM và danh mục}\\ khách hàng, chiến dịch,\\ sản phẩm};
  \node[mpbox=40mm, below=of crm] (email) {\textbf{Email}\\ gửi, mở, click,\\ hủy đăng ký};
  \begin{scope}[on background layer]
    \node[mpzone, fit=(ads)(email)] (src) {};
  \end{scope}
  \node[mpzonelabel, anchor=south west] at (src.north west) {Nguồn dữ liệu};

  % Cột 2: giải pháp
  \node[mphead=58mm, right=14mm of ads.north east, anchor=north west] (core) {MarketPulse 360};
  \node[mpbox=58mm, below=of core] (stream) {\textbf{Nhánh luồng}\\ Kafka $\rightarrow$ Flink\\ số liệu trực tiếp trong vài giây};
  \node[mpbox=58mm, below=of stream] (batch) {\textbf{Nhánh lô}\\ Spark, Iceberg, Dagster\\ lịch sử Bronze -- Silver -- Gold};
  \node[mpbox=58mm, below=of batch] (cert) {\textbf{Kiểm chứng}\\ dbt tests và đối soát tài chính\\ chỉ công bố release đúng từng đồng};
  \node[mpbox=58mm, below=of cert] (serve) {\textbf{Phục vụ}\\ ClickHouse, PostgreSQL\\ số liệu đã kiểm chứng và trực tiếp};
  \begin{scope}[on background layer]
    \node[mpzone, fit=(core)(serve)] (mp) {};
  \end{scope}
  \node[mpzonelabel, anchor=south west] at (mp.north west) {Giải pháp Kỹ thuật Dữ liệu};

  % Cột 3: ứng dụng và người dùng
  \node[mphead=42mm, right=14mm of mp.east, anchor=west, yshift=14mm] (app) {Marketing\\ Command Center};
  \node[mpbox=42mm, below=of app] (ui) {FastAPI + Next.js\\ 10 không gian làm việc};
  \node[mpbox=42mm, below=9mm of ui] (users) {\textbf{Các bên liên quan}\\ CMO, Performance, CRM,\\ Analyst, Product,\\ Data Engineer, Privacy};

  \draw[mparrow] (src.east) -- (mp.west |- src.east);
  \draw[mparrow] (mp.east |- app.west) -- (app.west);
  \draw[mparrow] (ui.south) -- node[right, mpnote] {ra quyết định} (users.north);
\end{tikzpicture}
}
    \caption{Bối cảnh hệ thống và luồng giá trị của MarketPulse 360}
    \label{fig:system-context}
\end{figure}

\subsection{Hai yêu cầu đối lập: tốc độ và độ chính xác}
\label{sec:c1-hai-yeu-cau}

Bài toán đặt ra hai yêu cầu kéo hệ thống về hai hướng ngược nhau:

\begin{itemize}
    \item \textbf{Tốc độ (độ trễ thấp).} Khi một chiến dịch đột ngột tăng chi phí hoặc tỷ lệ chuyển đổi sụt giảm, người quản lý quảng cáo cần biết trong vài giây đến vài phút chứ không phải sáng hôm sau. Yêu cầu này đòi hỏi xử lý luồng có trạng thái: khử trùng lặp, gom phiên và cập nhật chỉ số chiến dịch ngay khi sự kiện đến.
    \item \textbf{Độ chính xác được kiểm chứng.} Số liệu dùng để phân bổ ngân sách và báo cáo tài chính phải nối đúng hành trình khách hàng trong cửa sổ 30 ngày, loại bỏ phiên bot, tôn trọng lựa chọn đồng thuận (\texttt{consent\_analytics}, \texttt{consent\_marketing}) và cân khớp doanh thu ròng đến từng đồng. Những phép tính này cần dữ liệu lịch sử đầy đủ, dữ liệu đến trễ và khả năng tính lại.
\end{itemize}

Không một công cụ đơn lẻ nào đáp ứng được cả hai. MarketPulse 360 giải quyết bằng cách tách hai loại số liệu một cách tường minh: số liệu \emph{trực tiếp} (tạm thời, cập nhật liên tục) và số liệu \emph{đã kiểm chứng} (chỉ được công bố sau khi vượt qua toàn bộ kiểm thử và đối soát). Hai loại số liệu không bao giờ bị cộng lẫn vào nhau trên giao diện.

\section{Các bên liên quan (Stakeholders)}
\label{sec:c1-stakeholders}

Từ quy trình ra quyết định Marketing của một doanh nghiệp bán lẻ, Nhóm xác định 7 bên liên quan. Mỗi bên có yêu cầu về dữ liệu ở độ chi tiết, độ trễ khác nhau.

\subsection{Giám đốc Marketing (CMO)}
\label{sec:c1-sh-cmo}
Chịu trách nhiệm phân bổ tổng ngân sách giữa các kênh. CMO cần biết kênh nào tạo ra doanh thu ròng thực sự (sau hoàn tiền), tiến độ giải ngân so với kế hoạch và kết quả dự kiến khi chuyển ngân sách giữa các chiến dịch. Dữ liệu cần ở mức tổng hợp theo ngày, phải được kiểm chứng và có thể giải thích.

\subsection{Quản lý quảng cáo (Performance Marketing)}
\label{sec:c1-sh-performance}
Vận hành chiến dịch hằng ngày. Người dùng này cần chỉ số chiến dịch gần thời gian thực (phiên, đơn hàng, doanh thu theo chiến dịch), cảnh báo khi chi phí tăng đột biến hoặc vượt ngân sách, kèm bằng chứng để quyết định tạm dừng hay điều chỉnh giá thầu.

\subsection{Chuyên viên CRM (Lifecycle Marketing)}
\label{sec:c1-sh-crm}
Cần hồ sơ khách hàng hợp nhất (Customer 360), phân đoạn theo hành vi mua (RFM), mức độ tương tác email và danh sách khách hàng có cơ hội mua thêm hoặc có nguy cơ rời bỏ. Yêu cầu quan trọng là chỉ được kích hoạt khách hàng đã đồng ý nhận marketing.

\subsection{Chuyên viên phân tích dữ liệu (Data Analyst)}
\label{sec:c1-sh-analyst}
Cần so sánh các mô hình phân bổ doanh thu (First-Touch, Last-Touch, Linear), phân tích phễu chuyển đổi và cohort giữ chân khách hàng trên dữ liệu đã kiểm chứng, đồng thời truy vết được hành trình của từng đơn hàng khi số liệu bất thường.

\subsection{Quản lý thương mại điện tử (E-commerce)}
\label{sec:c1-sh-product}
Quan tâm đến hiệu quả ngành hàng, thiết bị mua sắm, giá trị đơn hàng, nhịp mua theo khung giờ và vị trí rơi rụng trong phễu từ xem sản phẩm đến thanh toán.

\subsection{Kỹ sư dữ liệu (Data Engineer)}
\label{sec:c1-sh-engineer}
Chịu trách nhiệm cho tính đúng và tính sẵn sàng của toàn bộ luồng dữ liệu: theo dõi sức khỏe hệ thống, kết quả kiểm thử chất lượng, nguồn gốc (lineage) của từng bản phát hành, khả năng phát lại và tính lại lịch sử khi công thức thay đổi mà không mất dữ liệu cũ.

\subsection{Chủ sở hữu dữ liệu (Data Owner)}
\label{sec:c1-sh-privacy}
Bảo đảm dữ liệu cá nhân chỉ được dùng đúng mục đích: lọc theo đồng thuận, ẩn danh hóa định danh trên giao diện, phân quyền theo vai trò và lưu vết kiểm toán mọi truy cập.

\section{Yêu cầu nghiệp vụ (Business requirements)}
\label{sec:c1-business-requirements}

Nhóm xác định \textbf{10 yêu cầu nghiệp vụ} (\texttt{BR-01} đến \texttt{BR-10}), chia thành sáu nhóm nhu cầu. Với mỗi nhóm nhu cầu, phần dưới đây nêu rõ dữ liệu cần có và lý do không thể đáp ứng nếu thiếu giải pháp Kỹ thuật Dữ liệu.

\subsection{Nhóm 1: Chiến dịch mang lại doanh thu (BR-01)}
\label{sec:c1-br-group1}
\textbf{BR-01: So sánh các mô hình phân bổ doanh thu.} Phân bổ doanh thu gộp, hoàn tiền và doanh thu ròng cho chiến dịch theo ba mô hình First-Touch, Last-Touch và Linear trên cửa sổ 30 ngày trước đơn hàng, từ đó tính ROAS theo từng mô hình. \emph{Vai trò của Kỹ thuật Dữ liệu:} phải nối click quảng cáo trả phí (trong luồng web) với đơn hàng (trong luồng đơn hàng) của cùng một hồ sơ khách hàng qua nhiều ngày, đồng thời bảo đảm tổng tiền được phân bổ bằng đúng tổng doanh thu, điều mà báo cáo riêng của từng nền tảng quảng cáo không làm được.

\subsection{Nhóm 2: Phản ứng tức thời (BR-02 -- BR-05)}
\label{sec:c1-br-group2}
\begin{itemize}
    \item \textbf{BR-02: Giám sát trực tiếp.} Hiển thị chỉ số chiến dịch (phiên theo từng bước phễu, đơn hàng, doanh thu) và lượng truy cập cửa hàng (Live View) trong vài giây sau khi sự kiện xảy ra.
    \item \textbf{BR-03: Cảnh báo có giải thích.} Phát hiện vượt ngân sách, chi phí tăng đột biến hoặc tỷ trọng hoàn tiền bất thường, mỗi cảnh báo kèm bằng chứng (ngưỡng, giá trị quan sát, phiên bản dữ liệu đã dùng).
    \item \textbf{BR-04: Theo dõi tiến độ giải ngân.} So sánh chi phí lũy kế với ngân sách cấu hình cho từng chiến dịch trong kỳ.
    \item \textbf{BR-05: Mô phỏng phân bổ lại ngân sách.} Ước tính doanh thu ròng khi phân bổ lại ngân sách dựa trên ROAS ròng lịch sử, công khai giả định tuyến tính của phép mô phỏng.
\end{itemize}
\emph{Vai trò của Kỹ thuật Dữ liệu:} xử lý luồng có trạng thái để có số liệu trong vài giây, nhưng cảnh báo và mô phỏng phải dựa trên số liệu chi phí đã kiểm chứng để không ra quyết định sai từ dữ liệu trùng lặp hoặc đến trễ.

\subsection{Nhóm 3: Chính xác đến từng đồng (BR-06)}
\label{sec:c1-br-group3}
\textbf{BR-06: Kiểm chứng trước khi công bố.} Mọi bản phát hành số liệu phải vượt qua kiểm thử dữ liệu tự động (dbt) và đối soát tài chính trước khi người dùng nhìn thấy; số liệu trực tiếp và số liệu đã kiểm chứng được tách biệt. \emph{Vai trò của Kỹ thuật Dữ liệu:} cần một quy trình công bố nguyên tử (atomic publication) và sổ cái vận hành để một bản phát hành lỗi không bao giờ thay thế bản đang phục vụ.

\subsection{Nhóm 4: Hành vi khách hàng (BR-07 -- BR-09)}
\label{sec:c1-br-group4}
\begin{itemize}
    \item \textbf{BR-07: Phễu chuyển đổi.} Đếm số phiên đạt từng bước \texttt{page\_view}, \texttt{view\_item}, \texttt{add\_to\_cart}, \texttt{begin\_checkout}, \texttt{purchase} theo ngày và chiến dịch.
    \item \textbf{BR-08: Hồ sơ Customer 360.} Hồ sơ hợp nhất theo khách hàng, dòng thời gian hoạt động, phân đoạn RFM có lưu phiên bản, mức tương tác email và nhịp mua theo khung giờ.
    \item \textbf{BR-09: Cohort, sản phẩm và địa lý.} Tỷ lệ giữ chân theo cohort tháng, doanh thu theo ngành hàng và thiết bị, phân bố doanh thu theo 13 tỉnh/thành phố.
\end{itemize}
\emph{Vai trò của Kỹ thuật Dữ liệu:} hợp nhất định danh (khách đã đăng nhập, khách ẩn danh, phiên) trên nhiều nguồn và nhiều tháng dữ liệu; chỉ một kho dữ liệu lịch sử có quản trị mới tính được cohort và Customer 360 nhất quán.

\subsection{Nhóm 5 và 6: Quản trị dữ liệu (BR-10)}
\label{sec:c1-br-group5-6}
\textbf{BR-10: Quản trị dữ liệu.} Loại khỏi số liệu đã kiểm chứng mọi sự kiện không có đồng thuận phân tích; chỉ kích hoạt phân đoạn với khách hàng có đồng thuận marketing; ẩn danh hóa định danh trên giao diện; phân quyền ba vai trò \texttt{viewer}, \texttt{operator}, \texttt{admin} kèm nhật ký kiểm toán; và tính lại lịch sử (backfill) một cách cô lập từ các snapshot Bronze đã ghim khi công thức thay đổi. \emph{Vai trò của Kỹ thuật Dữ liệu:} quản trị phải được áp dụng ngay trong đường ống, không thể thêm sau ở tầng báo cáo.

\subsection{Tổng hợp yêu cầu nghiệp vụ}
\label{sec:c1-br-summary}

Bảng~\ref{tab:br-matrix} tổng hợp ánh xạ giữa sáu nhóm nhu cầu, mười yêu cầu nghiệp vụ, các bên liên quan chính và năng lực Kỹ thuật Dữ liệu cần thiết. Chương~\ref{chap:ung-dung} đối chiếu lại từng yêu cầu với màn hình, API và bảng dữ liệu đã hiện thực.

\begin{table}[H]
\centering
\caption{Ma trận ánh xạ nhóm nhu cầu, yêu cầu nghiệp vụ, bên liên quan và năng lực Kỹ thuật Dữ liệu}
\label{tab:br-matrix}
\small
\renewcommand{\arraystretch}{1.25}
\begin{tabularx}{\textwidth}{>{\raggedright\arraybackslash}p{2.9cm} >{\raggedright\arraybackslash}p{3.4cm} >{\raggedright\arraybackslash}p{2.6cm} >{\raggedright\arraybackslash}X}
\toprule
\textbf{Nhóm nhu cầu} & \textbf{Yêu cầu} & \textbf{Bên liên quan} & \textbf{Kỹ thuật Dữ liệu} \\
\midrule
1. Chiến dịch nào mang lại doanh thu & BR-01 Phân bổ đa mô hình & CMO và Data Analyst & Nối click--đơn hàng 30 ngày, phân bổ số nguyên bảo toàn tổng tiền \\
2. Phản ứng tức thời & BR-02 Giám sát trực tiếp; BR-03 Cảnh báo; BR-04 Giải ngân; BR-05 Mô phỏng & Data Engineer, Data Analyst và CMO & Xử lý luồng có trạng thái, khử trùng lặp, chỉ số chi phí đã kiểm chứng \\
3. Báo cáo đúng từng đồng & BR-06 Kiểm chứng trước công bố & Data Engineer và CMO & dbt tests, đối soát tài chính, công bố nguyên tử \\
4. Hành vi và vòng đời & BR-07 Phễu; BR-08 Customer 360; BR-09 Cohort, sản phẩm, địa lý & Lifecycle Marketing và Data Analyst & Hợp nhất định danh, lịch sử nhiều tháng, mô hình Gold \\
5--6. Quyền riêng tư, tính lại lịch sử & BR-10 Quản trị dữ liệu & Data Owner Data Engineer & Lọc consent trong đường ống, RBAC, snapshot Bronze bất biến \\
\bottomrule
\end{tabularx}
\end{table}
